\section{Related Work and Research Gap}
UAV-communications surveys identify line-of-sight dominance, three-dimensional mobility, interference, and topology as central design variables \cite{zeng2019sky,mozaffari2019tutorial}; NR-V2X literature documents standardized sidelink physical-layer and resource-allocation procedures \cite{garcia2021nrv2x,3gpp38300}. UAV-specific studies use sidelink for communication-aware scheduling and multi-hop forwarding \cite{mishra2022cu2x}, evaluate its suitability for internal swarm communication \cite{varonen2024sidelink}, and analyze challenging site-specific propagation \cite{giannakoulas2025uxu}. Complementary work derives analytical multi-UAV interference/outage models \cite{lau2023outage}, jointly optimizes trajectory, frequency, and routing \cite{alam2024jtfr}, or targets synchronization and Doppler-aware access \cite{zhang2026sync}.

Our question is narrower: rather than embedding interference control in a routing/learning optimizer, we isolate swarm size, desired-link topology, exact frequency-resource partitioning, and spatial desired/interferer advantage while retaining NR-aware TBS/LDPC and sourced BLER behavior. A matched-seed robustness layer then tests whether conclusions attributed to density survive changes in pairing, resource allocation, link adaptation, and noise accounting.
